\documentclass{article}
\usepackage[utf8]{inputenc}
\usepackage{hyperref}
\usepackage{listings}
\usepackage{xcolor}
\usepackage{tikz}
\usetikzlibrary{shapes.geometric, arrows, positioning}
\usepackage{booktabs}
\usepackage{amsmath}

\definecolor{codegreen}{rgb}{0,0.6,0}
\definecolor{codegray}{rgb}{0.5,0.5,0.5}
\definecolor{codepurple}{rgb}{0.58,0,0.82}
\definecolor{backcolour}{rgb}{0.95,0.95,0.92}

\lstdefinestyle{mystyle}{
    backgroundcolor=\color{backcolour},   
    commentstyle=\color{codegreen},
    keywordstyle=\color{magenta},
    numberstyle=\tiny\color{codegray},
    stringstyle=\color{codepurple},
    basicstyle=\ttfamily\footnotesize,
    breakatwhitespace=false,         
    breaklines=true,                 
    captionpos=b,                    
    keepspaces=true,                 
    numbers=left,                    
    numbersep=5pt,                  
    showspaces=false,                
    showstringspaces=false,
    showtabs=false,                  
    tabsize=2
}

\lstset{style=mystyle}

\title{Spinning Up Modern: Technical Modernization Report}
\author{Kiet Huynh}
\date{April 2026}

\begin{document}

\maketitle

\section{Core Architectural Evolution}
This report details the technical implementation of the modernization of OpenAI's Spinning Up repository. The primary goal was to transition from legacy, unoptimized software to a production-grade, GPU-accelerated research environment using the latest available library versions.

\subsection{Gymnasium Migration and Environment Standardization}
The project successfully transitioned from \texttt{gym} to \texttt{gymnasium} (v1.0+). 

\begin{table}[h]
\centering
\caption{API and Environment Standardization}
\begin{tabular}{@{}lll@{}}
\toprule
Feature & Implementation Detail & Version/API \\ \midrule
Step Logic & Decomposed Terminal/Truncated signals & \texttt{obs, rew, term, trunc, info} \\
Reset Logic & Information dictionary support & \texttt{obs, info} \\
CartPole & Standardized on stable version & \texttt{CartPole-v1} \\
HalfCheetah & Standardized on MuJoCo free release & \texttt{HalfCheetah-v5} \\ \bottomrule
\end{tabular}
\end{table}

\section{Deep Learning Backends}

\subsection{PyTorch 2.5+: Graph Mode and GPU Coordination}
PyTorch implementations have been optimized for modern hardware and software standards:
\begin{itemize}
    \item \textbf{Kernel Fusion:} Integrated \texttt{torch.compile()} for JIT-optimized execution.
    \item \textbf{Memory Management:} Fixed host-device transfer bottlenecks by standardizing on \texttt{.detach().cpu().numpy()} for all environment interactions.
    \item \textbf{Type Safety:} Implemented comprehensive Python type hints for all Actor-Critic architectures.
\end{itemize}

\subsection{TensorFlow 2.15: Eager Execution Migration}
A total rewrite of the TensorFlow backend removed all legacy \texttt{tf.compat.v1} dependencies.
\begin{itemize}
    \item \textbf{Functional Models:} All algorithms now use \texttt{tf.keras.Model} subclassing for better modularity.
    \item \textbf{Automatic Differentiation:} Replaced static graphs with \texttt{tf.GradientTape} orchestration.
\end{itemize}

\begin{lstlisting}[language=Python, caption=Modern TF2 Policy Update Pattern]
with tf.GradientTape() as tape:
    loss_pi, kl, ent = compute_loss_pi(obs, act, adv, logp_old)
grads = tape.gradient(loss_pi, ac.pi.trainable_variables)
pi_optimizer.apply_gradients(zip(grads, ac.pi.trainable_variables))
\end{lstlisting}

\section{Hardware Acceleration: Dynamic Library Resolution}
A critical technical achievement was solving the \textbf{cuDNN Version Conflict} present in modern pip-based NVIDIA distributions.

\subsection{Implementation Detail: Library Mapping}
TensorFlow 2.15 expects \texttt{libcudnn.so.8}, while modern \texttt{nvidia-*} packages provide \texttt{v9}. We implemented a dynamic initialization hook in \texttt{setup\_tf\_gpu()}:
\begin{enumerate}
    \item \textbf{Automated Discovery:} Recursive scanning of \texttt{site-packages} for valid NVIDIA shared objects.
    \item \textbf{Symbolic Resolution:} Real-time creation of versioned symlinks (e.g., \texttt{v9} $\rightarrow$ \texttt{v8}) in a localized runtime directory.
    \item \textbf{LD\_PRELOAD Injection:} Dynamic environment patching before the first backend import to ensure 100\% GPU detection.
\end{enumerate}

\section{Distributed Execution: MPI-GPU Synchronization}
Modernized MPI support ensures correct coordination with CUDA devices:
\begin{itemize}
    \item \textbf{Weight Sync:} Synchronized model parameters across processes using rank-0 broadcasting.
    \item \textbf{Gradient Aggregation:} Global gradient averaging via \texttt{mpi\_avg} before optimizer updates.
    \item \textbf{Resource Isolation:} Configured \texttt{allow\_growth} to prevent MPI ranks from conflicting on VRAM allocation.
\end{itemize}

\section{Verification Results}
The system passed a rigorous 100\% sweep of the core algorithm suite (VPG, PPO, DDPG, TD3, SAC, TRPO) with no initialization warnings and full hardware acceleration.

\section{Citation}
\begin{lstlisting}[language=TeX]
@misc{SpinningUpModern2026,
    author = {Kiet Huynh},
    title = {{Spinning Up in Deep RL (Modernized)}},
    year = {2026},
    publisher = {GitHub},
    howpublished = {\url{https://github.com/tkiethuynh/spinningup}}
}
\end{lstlisting}

\end{document}
